%! Author = julian
%! Date = 17.04.26


\section{Model Evaluation} \label{sec:model-evaluation}
After training, it was important to evaluate the performance of the models, as conclusions about the propensity of seizures could only be drawn from the \glspl{fp}, if the models were decent at predicting seizures.
\glspl{fp} from a random or worse-than-chance classifier would not carry the intended meaning.

For the ensemble, the input features were scaled again using the z-score normalizer, which was previously fit on the train data.
Then, the models were used to issue scores (0--1) for all existing segments, which are the model's prediction of how likely that segment is to be preictal.
The segment scores were then averaged per clip to get the clip scores.
These, rather than segment scores, are used for the further evaluations.

We used two types of metrics used to evaluate the performance: clip-based metrics and \glspl{ebm}.
Here, clip-based metrics refer to the standard measures used to evaluate models in machine learning, which use the confusion matrix of \glspl{tp}, \glspl{fp}, \glspl{fn}, and \glspl{tn} to compute metrics, such as precision, sensitivity, accuracy, and the \gls{roc auc}.
On the other hand, \glspl{ebm} are specific to time-series data and seizure prediction and quantify more clinically relevant properties.
Here, the (relative) \gls{tifw} \parencite{mormann_seizure_2007, eberlein_machine_2025} and the (relative) number of seizures predicted were used.

\subsection{Time in False Warning} \label{subsec:tifw}
Relative \gls{tifw} is the \gls{eb} equivalent to the \gls{fpr}, i.e., the inverse of specificity.
That is to say it is minimized if non-preictal (false) samples are correctly classified as such.
However, when a prediction is issued from a $\sim$\SI{10}{\minute} clip, it is not issued for the timespan of that clip, but for that clip's \gls{sph}, i.e., the timespan $\sim$\num{5}--\SI{65}{\minute} after the clip's end.
This means that summing the duration of \gls{fp} clips is not sufficient to calculate the \gls{tifw}.
The clip's \glspl{sph} and overlaps thereof must be taken into account.
When a system unnecessarily alarms the patient, it causes them distress and inconvenience, which is why it is essential to minimize false alarms. % todo source

To calculate \gls{tifw}, only valid, non-preictal clips are selected.
For all such clips that issue a warning (\glspl{fp}), the \glspl{sph} are selected.
Overlapping \glspl{sph} are then merged, so that shared time is not counted multiple times.
The total duration of these non-overlapping false-warning intervals is the absolute \gls{tifw}.
The relative \gls{tifw} is obtained by normalizing absolute \gls{tifw} by the maximum possible warning time, which is the sum of the non-overlapping \gls{sph} intervals from all non-preictal clips.
This yields a metric bounded between 0 and 1 which enables comparable interpretation across patients and is consistent with the guidelines in \textcite{mormann_seizure_2007}.

\subsection{Event-Based Sensitivity} \label{subsec:eb-sensitivity}
It is also essential to quantify the relative number of seizures successfully predicted, which is also referred to as \gls{eb} sensitivity.
For preictal clips, a positive prediction is considered a true warning, and a seizure is counted as predicted if it falls within the \gls{sph} of at least one positively predicted preictal clip.
This yields the number of predicted seizures, and \gls{eb} sensitivity is then computed as the ratio of predicted to total seizures.
This metric therefore reflects clinically relevant event detection performance at seizure level rather than clip level.

\subsection{Threshold Optimization} \label{subsec:threshold-optimization}
Whether a clip issues a warning depends on whether its score is greater than or equal to a classification threshold.
To optimize it, metrics are calculated for all possible thresholds, i.e., all unique model scores.
To take both the \gls{eb} specificity (1 - rel. \gls{tifw}) and the \gls{eb} sensitivity into account, the harmonic mean between them is calculated:
\begin{equation} \label{eq:harmonic-mean}
    H(a,b)=\frac{2ab}{a+b}
\end{equation}
\begin{equation} \label{eq:eb-score}
    S_{\mathrm{EB}} = H\!\left(\mathrm{Sensitivity}_{\mathrm{EB}},\, 1-\mathrm{rel.\ TIFW}\right)
\end{equation}
The threshold that maximizes this \gls{eb} score is set as the best threshold and used to determine binary clip predictions in further steps.

The metrics are calculated for the train and test set independently.
The results for the test set yield a more accurate evaluation of the performance, since it is unseen data.
However, the optimal threshold is determined for the test set individually, rather than using the train set's threshold.
This again cannot be considered purely prospective performance, but improves the performance, which is an advantage for the cycle extraction, and the priority here.
